% Respuesta revisada al Comentario 1 del Revisor 1.
% Reemplaza el bloque \begin{reply} ... \end{reply} del Comentario 1 en ResponsesReviewers.tex.
% Los valores marcados \XX salen de run_paper_experiments.m corrido con SimQuad1.exe
% (results/<fecha>/resumen.csv, tabla_metricas.tex, tabla_masa.tex). Las referencias
% a ecuaciones usan las etiquetas del manuscrito; eq:ajuste_theta, eq:yaw_frame,
% sec:discrete y tab:params son nuevas.

\begin{reply}
We thank the reviewer for this careful comparison between the manuscript and the released code. The reviewer is right: the previous code and the previous manuscript described different controllers, and the earlier simulations did not validate the adaptation law or Theorem~2. We have rewritten both so that they describe the same controller. All equations of Section~2 are now implemented in a single function, \texttt{controlador\_paso.m}, in which every line is tagged with its equation label; all parameters are read from one configuration file, \texttt{cfg\_paper.m}, which is saved with every result; and a correspondence table between equations and code lines is included in the repository (\texttt{EQUATION\_MAP.md}). Every figure and table was regenerated with this code. The discrepancies were resolved as follows.

\begin{enumerate}
\item \textbf{Saturated position feedback and payload-velocity term.} Both have been removed. The outer loop now implements Eq.~\eqref{eq:base_reference} exactly,
\begin{equation*}
\mathbf{v}_{\mathrm{base}}=\dot{\mathbf{x}}_d+\begin{bmatrix}\dot{\mathbf{h}}_h^\star\\0\end{bmatrix}+K_{\mathrm{traj}}\,\tilde{\mathbf{x}},
\end{equation*}
with a linear position term and no payload-velocity term. The feedforward $\dot{\mathbf{h}}_h^\star$ is new in the manuscript: since $\mathbf{x}=\mathbf{x}_l+\mathbf{h}$, the UAV must move at the payload reference velocity plus the rate of the quasi-static cable inclination. Without it, the cable-error dynamics are driven by $-\dot{\mathbf{h}}^\star_h$, which produces a periodic tracking error on curved references.

\item \textbf{Sign of the quasi-static cable reference.} The sign in the code was correct and the manuscript was wrong. With $\mathbf{h}=\mathbf{x}-\mathbf{x}_l$, the horizontal force balance on the payload gives, to first order, $\ddot{\mathbf{x}}_{l,h}=(g/l)\,\mathbf{h}_h$. Hence
\begin{equation*}
\mathbf{h}_h^\star=\frac{l}{g}\,\ddot{\mathbf{x}}_{d,h},\qquad \dot{\mathbf{h}}_h^\star=\frac{l}{g}\,\dddot{\mathbf{x}}_{d,h},
\end{equation*}
i.e., the UAV leads the payload. Eq.~\eqref{eq:href} and Assumption~\ref{as:smallangle} have been corrected. Both quantities are now computed analytically from the reference, not by numerical differentiation. With this feedforward and under perfect velocity tracking, the linearised horizontal error dynamics are homogeneous, so the reference itself no longer forces a tracking error. The residual error observed in the simulator is due to forces outside this model, chiefly the aerodynamic drag on the payload and cable, which the kinematic loop, having no integral action, does not reject.
% VERIFICAR en SimQuad: con cfg.planta.air_friction = 0 y sin arrastre de la carga el error
% permanente de E3 deberia caer casi a cero (en la planta sustituta: 0.18 m -> 4 mm).

\item \textbf{Filters, saturation and leaky integral.} The saturation of the position and swing-compensation terms and the leaky integrator have been removed; the sliding-surface integral is a pure integral. The filters on the payload tracking error, on the yaw error and on the kinematic command have also been removed. One first-order filter remains, on the numerical derivative of the velocity reference $\dot{\mathbf{v}}_d$ used in the regressor (time constant 50~ms, 2.5 control periods). The discretisation of the cable-rate estimator, the sampling period (taken from the simulator clock, nominally 20~ms) and a command limit $|u_i|\le u_{\max}=5$~m/s, with conditional integration of the super-twisting state, are declared in a new subsection, ``Discrete-time implementation'' (Section~\ref{sec:discrete}). The command limit is never reached in the reported runs (fraction of saturated samples: \XX).

\item \textbf{Channels of the inner compensator.} The manuscript now states that the compensator acts on the two horizontal velocity channels, expressed in the UAV yaw frame in which the velocity commands are applied and in which the model of Ref.~\cite{Lucas_ICUAS2014} is defined (new Eq.~\eqref{eq:yaw_frame}). The vertical and yaw channels receive the kinematic command directly. The model, the regressor and the stability analysis have been rewritten for $n=2$, with one regressor $\boldsymbol{\phi}_i\in\mathbb{R}^{1\times2}$ and one parameter vector $\boldsymbol{\theta}_i\in\mathbb{R}^{2}$ per channel, and the dimensions of every term are stated explicitly.

\item \textbf{Adapted quantities.} The reviewer is correct: the previous code adapted the super-twisting gains $K_1,K_2$, whereas the manuscript adapted the model parameters. The previous adaptation law was also not implementable, because it was written in terms of $\tilde{\boldsymbol{\theta}}$, which depends on the unknown true parameters. Both manuscript and code now adapt the model parameters
\begin{equation*}
\theta_{1,i}=\frac{1}{\kappa_{u,i}},\qquad \theta_{2,i}=\frac{\kappa_{v,i}}{\kappa_{u,i}},\qquad i\in\{x,y\},
\end{equation*}
with fixed super-twisting gains. With $\mathbf{u}_i=\boldsymbol{\phi}_i\hat{\boldsymbol{\theta}}_i-w_i$ and $\boldsymbol{\phi}_i=[\,a_i\ \ v_i\,]$, $a_i=\dot v_{d,i}-\beta\tilde v_i$, the sliding dynamics are $\dot s_i=\kappa_{u,i}\boldsymbol{\phi}_i\tilde{\boldsymbol{\theta}}_i-\kappa_{u,i}w_i+\delta_i$, $\tilde{\boldsymbol{\theta}}_i=\hat{\boldsymbol{\theta}}_i-\boldsymbol{\theta}_i$. Once $s_i\equiv0$, the integral state of the super-twisting algorithm satisfies $\nu_i=\boldsymbol{\phi}_i\tilde{\boldsymbol{\theta}}_i+\delta_i/\kappa_{u,i}$, i.e., it reconstructs the equivalent control. The adaptation law uses this signal, which is measured:
\begin{equation*}
\dot{\hat{\boldsymbol{\theta}}}_i=\mathrm{Proj}_{\Theta}\Big(-\Gamma\,\boldsymbol{\phi}_i^{\top}\,\mathrm{sat}_{\bar\nu}(\nu_i)-\sigma\big(\hat{\boldsymbol{\theta}}_i-\boldsymbol{\theta}_{0}\big)\Big),
\end{equation*}
active only after $t_{\mathrm{on}}$ and while $|s_i|\le s_{\mathrm{on}}$, i.e., once the sliding mode has been reached (Eq.~\eqref{eq:ajuste_theta}). It is a gradient estimator on the parameter error. The projection keeps $\hat{\kappa}_{u,i}>0$ and bounds $\hat{\boldsymbol{\theta}}_i$. Convergence to the true parameters requires persistent excitation of $\boldsymbol{\phi}_i$ and a zero-mean $\delta_i$, and the revised manuscript states where these conditions fail: on the circular reference the speed is constant, so $a_i\approx0$ in the yaw frame and $\theta_{1,i}$ is not excited; the lateral channel is also biased by the centripetal term contained in $\delta_y$.
% VERIFICAR con E4/E7 en SimQuad (figuras *_5_theta.png): theta1 casi constante en el circulo,
% theta2_x hacia ~1, theta2_y sesgado. Si no se observa, reescribir esta frase.
Theorem~2 has been rewritten for this law (see our response to Comment~3).

\item \textbf{Nominal parameters and mismatched initialisation.} The previous text reported $\hat{\kappa}_u=0.5$ and $\hat{\kappa}_v=0.75$, whereas the code used different values. All values actually used are now listed in Table~\ref{tab:params}, which is generated from the configuration file. The main results use the nominal value implied by the velocity-command interface (unit static gain, $\kappa_u=\kappa_v=1/\tau$, $\tau=0.4$~s), i.e., $\boldsymbol{\theta}_0=[0.4,\ 1.0]^\top$. The former ``arbitrary'' values correspond to $\boldsymbol{\theta}=[2.0,\ 1.5]^\top$. With the model-inversion law they cannot serve as an initial condition: $\hat\theta_{2}>\theta_2$ introduces positive velocity feedback, the command saturates, and the sliding mode required to activate the adaptation is never reached. This is now stated in the manuscript. A new experiment (Section~\XX, Table~\XX) uses the lemniscate reference, which excites $\theta_{1}$, starts from $\hat{\boldsymbol{\theta}}_i(0)=[0.2,\ 0.5]^\top$, and compares fixed and adaptive parameters for payload masses of 0.05, 0.10, 0.20 and 0.30~kg. With adaptation, $\hat\theta_{2,x}$ converges to \XX, $\hat\theta_{1,x}$ increases with the payload mass from \XX\ to \XX, and the RMS of the integral state $\nu_x$ decreases from \XX\ to \XX, i.e., the identified model accounts for part of the effort previously supplied by the robust term. The payload tracking error is essentially unchanged (\XX\ vs.\ \XX~m), because the super-twisting algorithm already rejects this matched uncertainty; the manuscript no longer claims that adaptation improves tracking.

\item \textbf{Reproducibility.} A single script, \texttt{run\_paper\_experiments.m}, regenerates every figure and table without manual reconfiguration. The co-simulation runs in real time and is not strictly deterministic, so each case is repeated five times and reported as mean $\pm$ standard deviation. Payload velocities and all swing metrics are now computed from the simulator ground truth (the payload position returned by the simulator, differentiated with the simulator time stamps), not from the cable-rate estimate; the estimate itself is validated against this ground truth (see our response to Comment~4). The script also checks every analytical reference derivative against numerical differentiation.
\end{enumerate}
\end{reply}
